\documentclass[../../main.tex]{subfiles}

\begin{document}

\chapter{State of the Art} 
\label{chap:SOA}
Software visualization requires deciding \emph{how} software entities should be represented and, after deciding how to place them, how to display their connections. Different approaches have been proposed to address these problems, so in this chapter we first examine software visualization techniques in \secref{sec:soa_soft_vis}. We then introduce Wave Function Collapse in \secref{sec:soa_wfc}, explain how it works, and discuss its variants and applications in procedural content generation and constraint satisfaction problems in \secref{sec:wfc_variants}.

\section{Software Visualization Techniques}
\label{sec:soa_soft_vis} 

\subfile{graph-based.tex}
\subfile{polymetric_views.tex}
\subfile{3DandVR.tex}
\subfile{cartographic.tex}

\subfile{wfc-algorithm.tex}
\subfile{wfc-variants.tex}

\section{Summary}

In this chapter, we reviewed software visualization techniques, the standard Wave Function Collapse algorithm, and its applications and variants.

\paragraph{State of the Art.}
Software visualization uses different techniques to show software entities and their relationships. Graph-based layouts focus on positioning entities and drawing connections~\citebibliography{DiBa1998a, Fruc1991a, Tama1987a, Holt2006a}, while polymetric views add information about software metrics~\citebibliography{Lanz2003a, Agou2022a}. 3D and VR visualizations use spatial representations~\citebibliography{Wett2007a, Case2011a, More2023a, Hoff2024a, Gian2024a}, and cartographic approaches place entities based on their similarities~\citebibliography{Kuhn2010a}. WFC, originally developed for procedural content generation~\citeonline{Gumin2016wfc}, has also been extended in several ways. These extensions include global constraints~\citebibliography{Sandh2019a}, backtracking~\citebibliography{Karth2022a}, graph-based structures~\citebibliography{Kim2019a, Kim2025a}, and tiles of different sizes~\citebibliography{Piep2025a}. WFC is mainly used for generating game environments and levels~\citeonline{ParkerProcSkater, StalbergBadNorth, StalbergTownscaper, KleinebergInfiniteCity}.

\paragraph{Research Gap.}
Although both fields have been widely studied, their combination for software visualization remains limited. Existing visualization techniques use different methods to position entities and draw their connections. However, arranging paths while avoiding crossings and preserving all software relationships remains a challenge. WFC provides local compatibility rules and strategies for assigning tiles, but it does not directly address the specific requirements of software visualization. In particular, paths must connect the correct entities and preserve the relationships in the software system.

\paragraph{Influence of Previous Work.}
\wassis explores the use of WFC for drawing software relationships on a grid. It builds on existing techniques in both fields:
\begin{itemize}
    \item \textbf{Techniques from Previous Work:} From existing WFC extensions, we adopt dynamic tile weight adjustment~\citebibliography{Sandh2019a}, different backtracking strategies~\citebibliography{Karth2022a} and non-uniform tiles~\citebibliography{Piep2025a}.

    \item \textbf{Adaptations to Software Visualization:} We define tiles and compatibility rules based on the software relationships that must be displayed. We also combine WFC with MDS and introduce mechanisms to guide paths towards their destinations and handle contradictions. This approach allows us to investigate how WFC can be adapted to draw software relationships while preserving connectivity and reducing path crossings.
\end{itemize}

\end{document}